\documentclass[11pt,a4paper]{article}

\usepackage{fontspec}
\usepackage{babel}
\babelprovide[import, main]{russian}
\babelfont{rm}{Noto Serif}
\babelfont{sf}{Noto Sans}
\babelfont{tt}[AutoFakeBold=1.5]{Noto Mono}
\emergencystretch=2em

\usepackage[a4paper,margin=22mm,headheight=14pt]{geometry}
\usepackage{amsmath,amssymb,amsthm,mathtools}
\usepackage{booktabs,tabularx}
\usepackage{graphicx}
\usepackage{xcolor}
\usepackage{fancyhdr}
\usepackage{listings}
\usepackage{microtype}
\usepackage[hidelinks]{hyperref}

\definecolor{primary}{RGB}{24,62,108}
\definecolor{accent}{RGB}{138,43,43}
\definecolor{muted}{RGB}{85,95,108}
\definecolor{bgbox}{RGB}{244,247,251}
\definecolor{codebg}{RGB}{248,249,250}

\hypersetup{
  colorlinks=true,
  linkcolor=primary,
  citecolor=primary,
  urlcolor=primary,
  pdftitle={Мета-синсет и инвариант: типизированная система ревизии исследовательских представлений},
  pdfauthor={Александр Странд}
}

\setlength{\parskip}{5pt plus 2pt minus 1pt}
\setlength{\parindent}{0pt}

\pagestyle{fancy}
\fancyhf{}
\fancyhead[L]{\small\color{muted}Мета-синсет и инвариант: типизированная ревизия представлений (v2.4)}
\fancyhead[R]{\small\color{muted}А.\,Странд}
\fancyfoot[C]{\small\thepage}
\renewcommand{\headrulewidth}{0.4pt}
\renewcommand{\footrulewidth}{0pt}

\newtheoremstyle{preprintthm}
  {8pt}{6pt}{\itshape}{}{\bfseries\color{primary}}{.}{0.5em}{}
\newtheoremstyle{preprintdef}
  {8pt}{6pt}{\normalfont}{}{\bfseries\color{primary}}{.}{0.5em}{}

\theoremstyle{preprintthm}
\newtheorem{theorem}{Теорема}[section]
\newtheorem{lemma}[theorem]{Лемма}
\newtheorem{proposition}[theorem]{Предложение}

\theoremstyle{preprintdef}
\newtheorem{definition}[theorem]{Определение}
\newtheorem{remark}[theorem]{Замечание}
\newtheorem{axiom}{Аксиома}

\lstset{
  basicstyle=\ttfamily\footnotesize,
  backgroundcolor=\color{codebg},
  frame=single,
  rulecolor=\color{primary!25},
  xleftmargin=4pt,
  xrightmargin=4pt,
  aboveskip=6pt,
  belowskip=6pt,
  breaklines=true,
  keepspaces=true,
  columns=fullflexible,
  extendedchars=true,
  literate=
    {->}{{$\to$}}2
    {=>}{{$\Rightarrow$}}2
    {<=}{{$\le$}}2
    {>=}{{$\ge$}}2
    {forall}{{$\forall$}}1
    {exists}{{$\exists$}}1
    {Delta}{{$\Delta$}}1
    {Phi}{{$\Phi$}}1
    {lambda}{{$\lambda$}}1
    {kappa}{{$\kappa$}}1
    {pi}{{$\pi$}}1
    {neg}{{$\neg$}}1
    {lor}{{$\lor$}}1
    {land}{{$\land$}}1
}

\newcommand{\takeaway}[2]{%
  \vspace{4pt}\noindent\colorbox{bgbox}{%
    \parbox{\dimexpr\linewidth-2\fboxsep\relax}{%
      \vspace{3pt}\textbf{\color{primary}#1.}\hspace{0.4em}#2\vspace{3pt}%
    }%
  }\vspace{4pt}%
}

\newcommand{\LH}{\mathcal{L}_H}
\newcommand{\calB}{\mathcal{B}}
\newcommand{\calV}{\mathcal{V}}
\newcommand{\calR}{\mathcal{R}}
\newcommand{\calD}{\mathcal{D}}
\newcommand{\calU}{\mathcal{U}}
\newcommand{\Gain}{\operatorname{Gain}}
\newcommand{\Agree}{\operatorname{Agree}}
\newcommand{\Ent}{\operatorname{Ent}}
\newcommand{\dom}{\operatorname{dom}}
\newcommand{\im}{\operatorname{im}}

\begin{document}

\thispagestyle{plain}
\begin{center}
  {\LARGE\bfseries Мета-синсет и инвариант:\\[4pt]
  типизированная система ревизии исследовательских представлений}\\[10pt]
  {\large\color{muted} Дуальность экспансии и сужения, алгебра остаточного разрыва $\Delta(T)$, зависимая лицензия вывода и машинно верифицированное ядро в Lean~4.36.0}\\[12pt]
  {\normalsize Александр Странд}\\[4pt]
  {\small\color{muted} 23 ноября 2026\,г.\,$\cdot$\,препринт v2.4\,$\cdot$\,рукопись с верифицированным ядром v2.4 и Осью~2 (Gate~14--16 v2.4) в Lean~4.36.0}
\end{center}

\vspace{4pt}
\begin{center}
\begin{minipage}{0.92\linewidth}
  \begin{center}\textbf{Аннотация}\end{center}
  \small
  Исследовательский поиск распадается на два несовместимых по требованиям режима: порождение новых кандидатов за счёт междоменных аналогий требует управляемого дрейфа понятий, тогда как доказательная проверка и выбор следующего эксперимента требуют ссылочной прозрачности и запрета переписывать определения задним числом. В работе формализована дуальная архитектура двух контуров языка гипотез $\LH$ --- \texttt{meta-synset} (фокусная экспансия) и \texttt{invariant} (строгое сужение), --- связанных фазовым переходом \texttt{Freeze}, типизированным переносом \texttt{CertifiedTransfer} с обязательным остаточным разрывом $\Delta(T)=R\setminus S_{\mathrm{ok}}$ и ветвящимся шагом \texttt{StepBranch}. По сравнению с версией v1.0 устранены структурные дефекты наивной индексации областей: (1)~свидетельство \texttt{DirectEvidence s A c} жёстко привязано к утверждению $c:\texttt{Claim } s$ через функтор $\texttt{Denote}:\texttt{Claim } s\to\texttt{Prop}$; (2)~покрытое и полное обоснования индексированы \emph{точечно-зависимым} семейством $c:Q\to\texttt{Claim } s$, благодаря чему \texttt{evidenceAt} строится поточечно без гипотезы непустоты, а прямой обход моста полноты заблокирован теоремой \texttt{pointwise\_\allowbreak unproved\_\allowbreak blocks\_\allowbreak any\_\allowbreak total}; (3)~межобластной подъём \texttt{liftScope} использует условный \texttt{ScopeBridge} (действующий строго на \texttt{ModelOk}) и явно оставляет модельный разрыв \texttt{ModelGap}. Введена зависимая лицензия вывода $\kappa:\texttt{Judgement}(\alpha,\delta)$ с теоремами звучности, сужения языка (\texttt{restriction\_\allowbreak sound}) и анти-утечки, а для зонда ветвления \texttt{StepBranch} доказаны достаточность полного \texttt{SymbolLift} и необходимость каждого из двух условий (\texttt{binPreserving} и \texttt{injectiveBin}) по отдельности. Дополнительно формализована и верифицирована вторая ось эскалаций --- Ось~2 (Gate~14--16): эскалация наблюдательного порта по калибровочному сертификату \texttt{GaugeEquivCert} и запрет тупиковых скрытых переменных (Gate~14), алгебра вмешательств с вычислимой маской контроля $K$ и коммутатором $[T_L,T_R]$ под нормализатором $N$ (Gate~15), а также количественный остаточный разрыв $\Delta_\varepsilon$ (\texttt{HasGapBound}) и сертификаты произведения двух независимых e-value (фиксированный горизонт~2) на точных дискретных мерах (Gate~16). Все 44 ключевые декларации по четырём модулям (ядро v2.4 и Gate~14--16~v2.4) машинно верифицированы в едином прогоне Lean~4.36.0 (\texttt{0~errors}, \texttt{0~warnings}, \texttt{0~sorryAx}; 23 из 44 без аксиом, остальные используют лишь \texttt{propext}, \texttt{Classical.choice}, \texttt{Quot.sound}); дополнительно 12~теорем Gate~11--12~v12.3 верифицированы в~\cite{strand2026hl}.
\end{minipage}
\end{center}

\section{Проблема двух несовместимых режимов научного поиска}

Формальная теория идентификации и доказательной смены представления~\cite{strand2026hl,angluin1987} работает внутри уже заданной структуры: при фиксированном объектном языке $L=(G,J_G)$, наборе наблюдений $D$ и семантическом пространстве версий
\begin{equation}
  \calV(L,D) = \{f \in \calB(L) : \forall (x,y)\in D,\ f(x)=y\}, \qquad \calB(L)=\im J_G,
\end{equation}
можно строго определить согласованность, единственность идентификации $|\calV(L,D)|=1$, множество единогласия $\Agree(\calV)$ и одношаговое исключение языка $\calV(L,D\cup\{(q,y)\})=\varnothing$. Однако на этапе, когда старый язык уже опровергнут или понятийная рамка ещё не сложилась, исследователь сталкивается с двумя противоречащими требованиями:

\begin{enumerate}
  \item \textbf{Требование экспансии.} Чтобы разрушить понятийную фиксацию и найти нетривиальный класс-замену (например, перейти от счётчиковой памяти к стеку или очереди), необходимо сопоставлять объекты из далёких предметных областей под намеренно неполной предварительной линзой $\lambda_t$, допуская сдвиг и расщепление границ понятия.
  \item \textbf{Требование строгого сужения.} Чтобы сравнение кандидатов, отбраковка альтернатив и выбор следующего эксперимента не вырождались в подгонку определений задним числом, каждое утверждение обязано быть заморожено в своей версии, области и предпосылках, а любое свидетельство --- отделено от желаемого, но ещё не доказанного обобщения.
\end{enumerate}

Смешение двух режимов порождает две симметричные патологии: применение требований окончательного доказательства к этапу генерации запирает поиск в синонимах одного домена, а допуск плавающих определений на этап проверки порождает дегенеративный дрейф. Архитектура \texttt{meta-synset}\,$\leftrightarrow$\,\texttt{invariant} разводит эти режимы по двум изолированным контурам языка гипотез $\LH$ и соединяет их ровно одним типизированным швом --- фазовым переходом \texttt{Freeze} (рис.~\ref{fig:dual_loop}).

\begin{figure}[!htbp]
\centering
\begingroup\footnotesize
\begin{tabularx}{\linewidth}{@{}XcX@{}}
\toprule
Контур 1: \texttt{meta-synset} (фокусная экспансия) && Контур 2: \texttt{invariant} (строгое сужение)\\
\midrule
Семя $M_0=\langle\lambda_0,E_0\rangle$ && \texttt{Ranker}: четырёхзначное сравнение по $\Ent(T_2)\subseteq\Ent(T_1)$\\
$\texttt{Proj}(M_t,D)\to(E_D,\delta_D)$ &
$\xrightarrow{\ \texttt{Freeze}\ }$ &
\texttt{Research-Controller}: \texttt{StepBranch}, $\max_q\Gain_{\min}(q)$\\
$\texttt{Split}(M_t,\tau_{\mathrm{blind}}/\tau_{\mathrm{axis}})$: отсечение / ветвление &
$M_t\longmapsto\{\Phi_i\},\,\Delta(T)$ &
\texttt{CertifiedRev} / \texttt{ScopeBridge}: остаток $\Delta(T)$\\
\bottomrule
\end{tabularx}
\endgroup
\caption{Дуальная архитектура \texttt{meta-synset}\,$\leftrightarrow$\,\texttt{invariant}. Плавающий операнд $M_t$ не может напрямую поступить в \texttt{Ranker} без фазового перехода \texttt{Freeze}, уничтожающего дрейф и фиксирующего остаточный разрыв $\Delta(T)$.}
\label{fig:dual_loop}
\end{figure}

\section{Назначение \texorpdfstring{\texttt{meta-synset}}{meta-synset}: фокусная экспансия и управляемый дрейф}

В отличие от обычного синсета (WordNet), объединяющего синонимы внутри одного предметного слоя, \emph{мета-синсет} организует контролируемое накопление структурных аналогов из разных доменов и доводит их до точки бифуркации, в которой проступают скрытые оси различия~\cite{strand2026ms2}.

\begin{definition}[Мета-синсет и внутренний кортеж]
Пусть $\calU = \bigsqcup_{d\in\calD}\mathrm{Obj}_d$ --- размеченное объединение доменных объектов. Предварительная линза версии $t$ есть частичная процедура отбора $\lambda_t : \calU \rightharpoonup \mathcal{P}(\mathrm{Pred})$, обязанная уметь возвращать статус неприменимости без подбора новых признаков задним числом. \emph{Операнд-мета-синсет} версии $t$ имеет вид $M_t = \langle \lambda_t, E_t \rangle$, где $|\mathrm{domains}(M_t)| \ge 2$. Внутренний кортеж свидетельств внутри версии $M_t$ имеет координаты
\begin{equation}
  e = (d, I, s^-, s^+, u, K, \alpha, \delta, \kappa),
\end{equation}
где $I=(X,Y,Z,B)$ задаёт интерфейс и бюджет; $u:s^-\mapsto s^+$ --- допустимую операцию; $K$ --- маску удержанных, изменённых и неконтролируемых координат; $\alpha=(a,b,c,\eta)$ --- структурную карту; $\delta$ --- список утраченных и непроверенных свойств; $\kappa$ --- зависимую лицензию вывода (§~\ref{sec:judgement}).
\end{definition}

Все изменения операнда идут строго через типизированные операторы:
\begin{itemize}
  \item $\texttt{Proj}(M_t, D) \to (E_D, \delta_D)$: проекция линзы $\lambda_t$ в новый домен $D$ с анти-синонимическим фильтром и немедленной записью утраченных признаков $\delta_D$.
  \item $\texttt{Split}(M_t, \tau_{\mathrm{blind}}) \to M_{t+1}$: отсечение глухого контрпримера снаружи без расщепления операнда.
  \item $\texttt{Split}(M_t, \tau_{\mathrm{axis}}) \to (M_{t+1}^a, M_{t+1}^b)$: осевой тест, удерживающий одну скрытую ось и ломающий другую; фиксирует точку бифуркации, закрывая $M_t$ и порождая две новые ветви.
  \item $\texttt{Merge}(M_1, M_2) \to M_{12}$: слияние только при совместимых линзах и совпадении выделяемых структур на общих доменах.
  \item $\texttt{Freeze}(M_t) \to \{\Phi_1,\dots,\Phi_k\}$: фазовый переход, ликвидирующий дрейф и выпускающий замороженных кандидатов $\Phi=(P,S,\mathcal{A},\bot_\Phi)$.
\end{itemize}

Корректность экспансии удерживается тремя аксиомами --- \emph{ссылочной прозрачностью} (статус версии $M_{t+1}$ не переносится на $M_t$), \emph{анти-вырождением} ($|\mathrm{domains}(M_t)|\ge 2$) и \emph{завершением на бифуркации} ($M_t$ закрыт после осевого разбиения) --- и двумя операционными разрывами: фазовым разрывом между $M_t$ и $\Phi$ и разрывом воспроизводимости наблюдателя. Например, семя $M_0=\langle\lambda_{\mathrm{filter}},\{\text{firewall},\ \text{клеточная мембрана},\ \text{таможня}\}\rangle$ после проекции в нейросети ($\delta$: «исчез физический барьер») и глухого контрпримера «кирпичная стена» расщепляется осевым тестом «марковское одеяло» на две независимые ветви --- $M_2^a$ (интерфейсное экранирование) и $M_2^b$ (избирательный правиловый фильтр), --- и лишь затем $\texttt{Freeze}(M_2^b)$ передаёт замороженных кандидатов в \texttt{invariant}. В текущем Lean-ядре контур \texttt{meta-synset} механизирован на базовом уровне (\texttt{MetaOperand}, \texttt{metaProj}, \texttt{metaFreeze}), тогда как полная семантика \texttt{Split}/\texttt{Merge} и семь осей \texttt{Ranker} остаются формальной спецификацией.

\section{Назначение \texorpdfstring{\texttt{invariant}}{invariant}: частичный порядок следствий и контроллер проб}

Контур \texttt{invariant} отвечает за два вопроса, недоступных для \texttt{meta-synset}: \emph{что именно установлено} (\texttt{Ranker}) и \emph{какой эксперимент провести следующим} (\texttt{Research-Controller}).

\paragraph{Разведение носителя и содержания.}
Носитель свидетельства (\texttt{Carrier}: ручной вывод, терм ядра Lean, конечный сертификат, исчерпывающий перебор, эмпирический протокол) не определяет автоматически силу вывода: один и тот же \texttt{LeanKernelTerm} может доказывать как частное свойство на четырёх конечных историях, так и универсальную теорему на $\Sigma^*$~\cite{strand2026ms2}. Поэтому \texttt{Ranker} оценивает кандидатов по семи \emph{осям вопросов} (Proof, Scope, Robustness, Causality, Transfer, Utility, Discrimination) без искусственного суммарного балла, используя предпорядок доминирования $T_1 \succeq T_2 \iff \Ent(T_2)\subseteq\Ent(T_1)$ и четырёхзначный исход сравнения: \emph{доминирует}, \emph{доминируем}, \emph{несравнимы} либо \emph{недостаточно сведений}. При этом \emph{положение утверждения} строго отделяется от \emph{очереди исследований}: слабая гипотеза с дешёвым решающим контрпримером стоит первой в очереди \texttt{Research-Controller}.

\section{Топология областей, запечатанные свидетельства и алгебра разрыва \texorpdfstring{$\Delta(T)$}{Delta(T)}}
\label{sec:topology}

Связующий объект \texttt{CertifiedTransfer} (раздел~1) реализуется в ядре через структуры \texttt{TransferCoverage} и \texttt{CoveredJustification}: остаточный разрыв $\Delta(T)=R\setminus S_{\mathrm{ok}}$ тождественен множеству $\{q\in R\mid \mathrm{Gap}(q)\}$, а условие $\Delta(T)=\varnothing$ --- предикату \texttt{GapEmpty Gap}.

В ранней схеме v1.0 индексация по одной лишь области \texttt{Scope} оставляла возможность подменять утверждения внутри области и скрывать непустой разрыв при подъёме \texttt{Lift}. В ядре v2.1 разделены \texttt{CoveredJustification} и \texttt{TotalJustification}, однако утверждение $c:\texttt{Claim } s$ оставалось константным по $Q$, из-за чего непустой разрыв блокировал только \texttt{bridgeComplete}, не запрещая собрать \texttt{TotalJustification} напрямую. В верифицированном ядре v2.4 (\texttt{epistemic\_\allowbreak revision\_\allowbreak kernel\_\allowbreak v2\_\allowbreak 4.lean}) эта граница закрыта точечной зависимостью $c : Q \to \texttt{Claim } s$.

\paragraph{Топология четырёх областей.}
Все утверждения разбиты на четыре непересекающиеся области \texttt{Scope}: \texttt{FinitePrefix} (конечный префикс строго добавляемого журнала \texttt{RawLog}), \texttt{LanguageMenu} (замороженное до опыта конечное меню альтернатив $\calR$), \texttt{InfiniteLanguage} (всё бесконечное множество слов $\Sigma^*$) и \texttt{RevisionSpace} (пространство трасс и слоёв ревизии $\Delta_L/\Delta_D/\Delta_\pi$).

\begin{lstlisting}[caption={Верифицированное ядро v2.4 (Lean~4.36.0): точечно-зависимые утверждения $c:Q\to\texttt{Claim } s$, поточечный \texttt{evidenceAt} и условный \texttt{ScopeBridge}.}]
structure CoveredJustification (Q : Type) (s : Scope) (A : Premise) (c : Q -> Claim s) where
  R          : Q -> Prop
  Covered    : Q -> Prop
  Gap        : Q -> Prop
  partition  : forall q : Q, R q -> Covered q lor Gap q
  localProof : A -> forall q : Q, Covered q -> Denote (c q)

structure TotalJustification (Q : Type) (s : Scope) (A : Premise) (c : Q -> Claim s) where
  R          : Q -> Prop
  evidenceAt : forall q : Q, R q -> DirectEvidence s A (c q)
  fullCover  : A -> forall q : Q, R q -> Denote (c q)

structure ScopeBridge (Q2 : Type) (s1 s2 : Scope) (A1 A2 : Premise)
    (c1 : Claim s1) (c2 : Q2 -> Claim s2) (ModelOk : Q2 -> Prop) where
  liftProof : A2 -> Denote c1 -> forall q : Q2, ModelOk q -> Denote (c2 q)
\end{lstlisting}

\begin{figure}[!htbp]
\centering
\begingroup\footnotesize
\begin{tabularx}{\linewidth}{@{}XcXX@{}}
\toprule
Covered (\texttt{FinitePrefix}) &
$\xrightarrow[\texttt{GapEmpty}]{\ \texttt{bridgeComplete}\ }$ &
Total (\texttt{FinitePrefix}) & \texttt{LanguageMenu}\\
$\texttt{partition}:R\Rightarrow\mathrm{Cov}\lor\mathrm{Gap}$ &&
\texttt{DirectEvidence s A c} &
Menu / Split; \texttt{StepBranch}, \texttt{SymbolLift}, $\Gain_{\min}$\\
\midrule
Covered (\texttt{InfiniteLanguage}) &
$\xleftarrow{\ \texttt{liftScope}\ }$ &
Подъём из Total (\texttt{FinitePrefix}) через \texttt{ScopeBridge} &
\texttt{RevisionSpace}\\
$w\in\Sigma^*$; \texttt{ModelGap} (\texttt{hasEpsilon}) &&
Свидетель разрыва $m$: $\mathrm{hasEps}(m)$ &
$\mathrm{Obs}$, $\Delta L/\Delta D/\Delta\pi$\\
\bottomrule
\end{tabularx}
\endgroup
\caption{Топология областей \texttt{Scope} и два различных моста ядра v2.4: мост полноты \texttt{bridgeComplete} (закрытие разрыва при \texttt{GapEmpty} внутри одной области) и межобластной подъём \texttt{liftScope} (перевод \texttt{Total} из \texttt{FinitePrefix} в \texttt{Covered} в \texttt{InfiniteLanguage} с явным модельным разрывом \texttt{ModelGap}).}
\label{fig:scope_lattice}
\end{figure}

Отметим асимметрию сигнатуры \texttt{ScopeBridge}: источник $c_1:\texttt{Claim } s_1$ одиночен, а цель $c_2:Q_2\to\texttt{Claim } s_2$ является точечным семейством, поскольку один конечный сертификат согласия на трассе (\texttt{TraceAgrees}) поднимается параметрической теоремой сразу на целое семейство целевых моделей $Q_2$.

\begin{theorem}[Закрытие лазеек и точечная алгебра разрывов; верифицировано в Lean~4.36.0]
\label{thm:kernel_core}
В модуле \texttt{EpistemicRevisionV2} (v2.4, без аксиом ядра) доказаны следующие свойства:
\begin{enumerate}
  \item \emph{Интенсиональность (\texttt{claim\_\allowbreak distinct\_\allowbreak not\_\allowbreak reusable}):} наличие свидетельства для \texttt{cTrue} не позволяет построить \texttt{DirectEvidence} для \texttt{cFalse} в той же области.
  \item \emph{Поточечный мост без гипотезы непустоты (\texttt{bridgeComplete}):} при \texttt{GapEmpty Gap} конструкция \texttt{evidenceAt q hRq} строится поточечно из \texttt{partition} и \texttt{localProof}, делая прежнюю гипотезу непустоты запроса излишней.
  \item \emph{Блокировка любого прямого тотала (\texttt{pointwise\_\allowbreak unproved\_\allowbreak blocks\_\allowbreak any\_\allowbreak total}):} если $\exists q_0\in Q,\ R(q_0)\land A\land \neg\texttt{Denote}(c\ q_0)$, то \texttt{TotalJustification} с точным равенством области \texttt{tot.R = R} \emph{не существует в принципе} (нельзя собрать даже в обход \texttt{bridgeComplete}; при этом тотал на суженной области $R'\subset R$, обходящей $q_0$, не исключается).
  \item \emph{Композиция разрывов (\texttt{composeCoverage\_\allowbreak gap\_\allowbreak empty\_\allowbreak iff}):} вдоль $f:Q_1\to Q_2$ составная покрытая область есть $S_{\mathrm{ok},1}(q)\land S_{\mathrm{ok},2}(f(q))$, составной разрыв --- $\mathrm{Gap}_1(q)\lor \mathrm{Gap}_2(f(q))$.
  \item \emph{Условный подъём (\texttt{liftScope}, \texttt{gate12\_\allowbreak 13\_\allowbreak epsilon\_\allowbreak blocks\_\allowbreak any\_\allowbreak total}):} \texttt{ScopeBridge} действует строго на подмножестве \texttt{ModelOk}; если существует модель $m_0$ с $\varepsilon$-переходами, в которой выполнены $A_1\land A_2$, но целевое утверждение ложно, то \texttt{TotalJustification} с \texttt{tot.R = (fun \_ => True)} на всём пространстве моделей $Q_{\mathrm{mod}}$ невозможен.
\end{enumerate}
\end{theorem}

\takeaway{Принципиальное различие двух мостов}{В ядре v2.4 \texttt{bridgeComplete} и \texttt{liftScope} --- разные математические операции. Первый закрывает остаток $\Delta(T)=\varnothing$ внутри фиксированной области; второй переносит результат между областями и \emph{меняет природу разрыва} (от конечной выборки слов к модельному сужению real-time \texttt{StackProgram} против общего PDA).}

\section{Зависимая лицензия вывода \texorpdfstring{$\kappa : \texttt{Judgement}(\alpha,\delta)$}{Judgement} и динамика \texttt{RawLog}}
\label{sec:judgement}

Вместо перечисления ярлыков методов координата $\kappa$ в кортеже мета-синсета формализована как тип лицензии на разрешённые выводы $\Gamma(\kappa)=\texttt{AllowedConclusion }\kappa$. Конструкторы включают \texttt{Transport}, \texttt{LocalIso}, \texttt{Refutation}, \texttt{NonIdent}, \texttt{Expansion}, \texttt{Restriction}, \texttt{FormalTheorem} и \texttt{FiniteCheck}.

Для точной формулировки теорем о лицензиях напомним используемые в ядре определения: $\mathrm{LangEmbed}(L_2, L_1) \iff \forall h_2,\ \exists h_1,\ \forall w,\ L_1.\mathrm{eval}(h_1,w)=L_2.\mathrm{eval}(h_2,w)$ и $\mathrm{VersionSpaceUnique}(L, D) \iff \forall h_1,h_2,\ (\mathrm{agreesOn}\ L\ D\ h_1 \land \mathrm{agreesOn}\ L\ D\ h_2 \to \forall w,\ L.\mathrm{eval}(h_1,w)=L.\mathrm{eval}(h_2,w))$ (т.е.\ $|\tilde V(L,D)|\le 1$, включая случай пустого пространства версий).

\begin{theorem}[Звучность сужения и теоремы анти-утечки; верифицировано в Lean~4.36.0]
\label{thm:judgement}
\begin{enumerate}
  \item \emph{Звучность сужения (\texttt{restriction\_\allowbreak sound}, 0 аксиом):} для любых языков с семантическим вложением $\mathrm{LangEmbed}(L_2, L_1)$ (каждая программа из $L_2$ эквивалентна некоторой программе из $L_1$) сужение языка сохраняет как опровержение ($\mathrm{VersionSpaceEmpty}\ L_1\ D\to\mathrm{VersionSpaceEmpty}\ L_2\ D$), так и семантическую единственность ($\mathrm{VersionSpaceUnique}\ L_1\ D\to\mathrm{VersionSpaceUnique}\ L_2\ D$).
  \item \emph{Анти-утечка конечной проверки (\texttt{finiteCheck\_\allowbreak blocks\_\allowbreak universal}, \texttt{propext}):} существуют граница $B$ и предикат $P$, для которых $\Gamma(\texttt{FiniteCheck } B\ P)$ истинно, но $\forall w\in\Sigma^*,\ P(w)=\mathrm{true}$ ложно.
  \item \emph{Анти-утечка опровержения (\texttt{refutation\_\allowbreak blocks\_\allowbreak extended}, \texttt{propext}):} существуют языки $L_1, L_2$ и журнал $D$, для которых $\Gamma(\texttt{Refutation } L_1\ D)$ истинно, но $L_2$ не опровергнут на $D$.
  \item \emph{Префиксная монотонность (\texttt{versionSpaceEmpty\_\allowbreak monotone}, 0 аксиом):} если язык $L$ опровергнут на префиксе $l_1$ строго добавляемого журнала (\texttt{PrefixOf $l_1$ $l_2$}), он остаётся опровергнутым на любом расширении $l_2$.
\end{enumerate}
\end{theorem}

Поверх префиксной монотонности задаются необратимо замороженное меню \texttt{FrozenMenu} (привязанное к префиксу \texttt{frozenAt}), удостоверенная ревизия \texttt{CertifiedRev} (состояние до/после, доказательство \texttt{PrefixOf}, список сырых редакций $\Delta_L/\Delta_D/\Delta_\pi$ и \texttt{CoveredJustification}) и сквозная цепочка \texttt{DiscoveryChain}. Вместе с теоремой о слоях ревизии~\cite[Теорема~3]{strand2026ms2} это фиксирует, что внешние ответы без журнала грамматики и данных не локализуют ревизию.

\section{Ветвящийся контроллер \texorpdfstring{\texttt{StepBranch}}{StepBranch}, \texorpdfstring{\texttt{SymbolLift}}{SymbolLift} и минимальный зонд \texorpdfstring{$\Gain_{\min}(q)$}{Gain\_min(q)}}

Переход от жёсткой линейной лестницы (\texttt{DFA}\,$\to$\,\texttt{1-DCA}\,$\to$\,\texttt{DPDA}) к дереву альтернатив задаётся шагом \texttt{StepBranch}: после опровержения текущего класса на префиксе $D_{\le t}$ контур \texttt{meta-synset} поставляет замороженное меню $\calR$, а \texttt{Research-Controller} выбирает зонд $q$ по гарантированному минимальному отсечению:
\begin{equation}
  \Gain_{\min}(q) = |\calR| - \max_{y} |\calR_{q,y}|.
\end{equation}

В аудиторском разборе v1.0 было выявлено, что для 3-точечного зонда разделения стека (LIFO) и очереди (FIFO) недостаточно знать лишь $f(0)=b_0$ и $f(1)=b_1$: необходим явный объект спецификации \texttt{SymbolLift f}, удостоверяющий посимвольное гомоморфное продолжение $f(u\mathbin{+\!\!+} v)=f(u)\mathbin{+\!\!+} f(v)$, сохранение бинарного алфавита (\texttt{binPreserving}: $f(0),f(1)\in\{0,1\}$) и инъективность (\texttt{injectiveBin}: $f(0)\ne f(1)$).

\begin{theorem}[Зависимость зонда от нулевого класса и необходимость \texttt{SymbolLift}]
\label{thm:probe_complexity}
\begin{enumerate}
  \item \emph{Rule~110 на бесконечной $\mathbb{Z}$~\cite[Теоремы~4--5]{strand2026ms2}:} против произвольной эффективной аффинной карты $A_t$ минимальный вес подготовки равен $2$ ($F_t(a,b)=a\lor b$ против $A_t(a,b)=a\oplus b$); при сужении нулевого класса до итерации одного трансляционно-инвариантного аффинного правила $U^t$ на чётном горизонте минимальный вес падает до $1$.
  \item \emph{Разделение LIFO/FIFO в \texttt{StepBranch} (\texttt{stepBranch\_\allowbreak order\_\allowbreak probe\_\allowbreak splits}, Lean~4.36.0):} при наличии \texttt{SymbolLift f} 5-символьный зонд порядка $q^*=[0,1,\#,f(1),f(0)]$ даёт \texttt{true} на канонической стековой ветви и \texttt{false} на канонической очередевой ветви ($\Gain_{\min}(q^*)=1$).
  \item \emph{Необходимость \texttt{binPreserving} (\texttt{stepBranch\_\allowbreak fails\_\allowbreak without\_\allowbreak binPreserving}, Lean~4.36.0):} для отображения $f_{\mathrm{bad}}(0)=\#,\ f_{\mathrm{bad}}(1)=0$ условие инъективности $f_{\mathrm{bad}}(0)\ne f_{\mathrm{bad}}(1)$ выполнено, однако обе ветви возвращают \texttt{false} на $q^*$ ($\Gain_{\min}(q^*)=0$).
  \item \emph{Необходимость \texttt{injectiveBin} (\texttt{stepBranch\_\allowbreak fails\_\allowbreak without\_\allowbreak injectiveBin}, Lean~4.36.0):} для отображения $f_{\mathrm{const}}(0)=f_{\mathrm{const}}(1)=0$ условие \texttt{binPreserving} выполнено, однако обе ветви возвращают \texttt{true} на $q^*$ ($\Gain_{\min}(q^*)=0$).
\end{enumerate}
\end{theorem}


\section{Ось~2: эскалация порта (Gate~14), алгебра вмешательств (Gate~15) и количественный остаток (Gate~16)}
\label{sec:axis2}

Иерархия памяти \texttt{DFA}\,$\to$\,\texttt{1-DCA}\,$\to$\,\texttt{DPDA}\,$\to$\,\texttt{Queue} (Ось~1) срабатывает, когда пространство версий \emph{пустеет} ($|\calV|=0$) на фиксированном порте наблюдений. Двойственная проблема возникает при анализе многокомпонентных систем и больших языковых моделей: ни один допустимый запрос на том же публичном порте не различает сертифицированную калибровочную пару, хотя другие кандидаты могут быть исключены. Это требует ортогональной Оси~2 --- эскалации самого наблюдательного интерфейса, алгебры вмешательств и перехода к вероятностным оценкам остаточного разрыва (рис.~\ref{fig:two_axis}).

\begin{figure}[!htbp]
\centering
\begingroup\footnotesize
\begin{tabularx}{\linewidth}{@{}XXX@{}}
\toprule
\multicolumn{3}{@{}l}{Ось 1: эскалация памяти при $|\calV|=0$ на фиксированном порте}\\
\midrule
Gate 11: \texttt{MNCert}; DFA($k$)$\to$1-DCA &
Gate 12 v12.3: \texttt{1DCACert}; $\forall m\,\neg$1-DCA; переход к DPDA &
Gate 13 (кандидат): \texttt{StepBranch} и Queue ($\forall m\,\neg$DPDA)\\
\midrule
\multicolumn{3}{@{}l}{Ядро v2.4: \texttt{Scope}, Covered/Total и $\Delta(T)$}\\
\midrule
\multicolumn{3}{@{}l}{Ось 2: порт, вмешательства и статистическая проверка при удостоверенном препятствии публичного порта или шуме}\\
\midrule
Gate 14: \texttt{GaugeEquivCert}, \texttt{tapPort}; запрет \texttt{SpuriousExtension} &
Gate 15: маска $K$ и коммутатор $[T_L,T_R]$ &
Gate 16: \texttt{HasGapBound}, $\Delta_\varepsilon$ и e-value, $\Pr(\text{отклонение})\le 1/K$\\
\bottomrule
\end{tabularx}
\endgroup
\caption{Двухосевая карта исследовательских эскалаций: Ось~1 (смена класса памяти при $|\calV|=0$) и верифицированная в Lean~4.36.0 Ось~2 (эскалация порта, алгебра вмешательств и количественный остаток при $|\calV|>1$ или шуме).}
\label{fig:two_axis}
\end{figure}

\begin{theorem}[Ось~2: порт, вмешательства и количественный остаток; верифицировано в Lean~4.36.0]
\label{thm:axis2}
В трёх независимых модулях чистого Lean~4 Core (без \texttt{sorry}) доказаны следующие результаты:
\begin{enumerate}
  \item \emph{Gate~14 (\texttt{gate14\_\allowbreak port\_\allowbreak escalation\_\allowbreak v2\_\allowbreak 4.lean}, 10 теорем~\cite{strand2026g14}):} для двухступенчатого механизма $Q\xrightarrow{R}Z\xrightarrow{S}Y$ и обратимой замены $T$ (\texttt{GaugeEquivCert}: $T(R(q_0))\ne R(q_0)$) калибровочный двойник $M_2=(T\circ R, S\circ T^{-1})$ отличается от $M_1$, но абсолютно неразличим на любом журнале через \texttt{publicPort} (\texttt{gaugeCert\_\allowbreak blocks\_\allowbreak single\_\allowbreak port\_\allowbreak id}). Уточнение порта \texttt{PortRefinement} (нестрогое условие $\mathrm{proj}\circ P_2 = P_1$) сохраняет прежние опровержения, а расширенный порт \texttt{tapPort} разбивает калибровочный слой за один запрос $q_0$ (\texttt{tapPort\_\allowbreak breaks\_\allowbreak gaugeCert}). Отдельно доказано, что тупиковое расширение \texttt{SpuriousExtension} с неизменным наблюдаемым образом не может починить опровергнутый класс ни на каком журнале (\texttt{spurious\_\allowbreak extension\_\allowbreak cannot\_\allowbreak repair\_\allowbreak refutation}).
  \item \emph{Gate~15 (\texttt{gate15\_\allowbreak intervention\_\allowbreak algebra\_\allowbreak v2\_\allowbreak 4.lean}, 8 теорем~\cite{strand2026g15}):} для двух вмешательств $(T_L, T_R)$ под формальным нормализатором $N$ локальная ортогональность маски $K$ в точке $x_0$ проверяется вычислимым предикатом \texttt{checkOrthogonalControl} (\texttt{by decide}). Если $T_L$ глобально сохраняет $N$, коммутирование $N(T_L(T_R x))=N(T_R(T_L x))$ равносильно языковой инвариантности стратегии $T_R$; один свидетель некоммутативности опровергает совместную инвариантность (\texttt{noncomm\_\allowbreak refutes\_\allowbreak joint\_\allowbreak invariance}). Теоремы анти-утечки (\texttt{confounded\_\allowbreak TL\_\allowbreak fakes\_\allowbreak noncomm\_\allowbreak without\_\allowbreak mask\_\allowbreak K}, \texttt{unaligned\_\allowbreak normalizer\_\allowbreak fakes\_\allowbreak semantic\_\allowbreak shift}) дают два разных контрпримера: нарушение удержания \texttt{Held} может вызвать зависимость от порядка даже при нечувствительной к языку стратегии, а невыровненный $N$ может создать ложный одношаговый семантический сдвиг даже при коммутирующих вмешательствах.
  \item \emph{Gate~16 (\texttt{gate16\_\allowbreak stochastic\_\allowbreak gap\_\allowbreak evalue\_\allowbreak v2\_\allowbreak 4.lean}, 11 деклараций~\cite{strand2026g16}):} на точных дискретных мерах в $\mathbb{N}$ доказано неравенство Маркова (\texttt{markov\_\allowbreak ineq}) и граница ложного отклонения для валидных e-переменных $\mathbb{E}_\mu[E/D]\le 1$: $K\cdot\mathrm{eventMass}(\mu,\{E\ge KD\})\le\mathrm{totalMass}(\mu)$ (\texttt{evalue\_\allowbreak false\_\allowbreak rejection\_\allowbreak bound}), а также мультипликативность произведения двух независимых e-переменных (\texttt{indep\_\allowbreak two\_\allowbreak step\_\allowbreak evalue\_\allowbreak product\_\allowbreak valid}). Количественный разрыв $\mathrm{HasGapBound}(\mu,\mathrm{Gap},\mathrm{num},\mathrm{den}) \iff (0 < \mathrm{den}) \land (\mathrm{eventMass}(\mu,\mathrm{Gap})\cdot\mathrm{den} \le \mathrm{num}\cdot\mathrm{totalMass}(\mu))$ исключает вакуумный нулевой знаменатель (\texttt{den\_\allowbreak zero\_\allowbreak rejected\_\allowbreak by\_\allowbreak HasGapBound}, 0~аксиом), обобщает \texttt{GapEmpty} (\texttt{gapEmpty\_\allowbreak implies\_\allowbreak zero\_\allowbreak gapBound}) и удовлетворяет субаддитивности при общем знаменателе (\texttt{union\_\allowbreak gap\_\allowbreak bound}).
\end{enumerate}
\end{theorem}


\paragraph{Операторы \texttt{meta-synset} и условная развилка Оси~1\,$\leftrightarrow$\,Оси~2.}
В ядре v2.4 добавлены базовые структуры \texttt{MetaOperand} (с инвариантом анти-вырождения $2\le\mathrm{domainCount}$), \texttt{metaProj} и \texttt{metaFreeze} (\texttt{metaFreeze\_\allowbreak preserves\_\allowbreak candidates}), а также условная развилка \texttt{chooseConditionalTwoAxisAction} / \texttt{conditional\_\allowbreak two\_\allowbreak axis\_\allowbreak routing\_\allowbreak blocks\_\allowbreak spurious\_\allowbreak memory\_\allowbreak escalation}: при постулированной несовместимости причин (\texttt{disjointCause}) калибровочная неидентифицируемость (\texttt{axis2Gauge}) блокирует слепую эскалацию памяти по Оси~1. Подчеркнём границу: \texttt{axis1Refuted} и \texttt{axis2Gauge} здесь являются свободными предикатами (ещё не выведенными внутри единого типа из \texttt{VersionSpaceEmpty} и \texttt{GaugeEquivCert}), поэтому данная теорема задаёт именно \emph{условную} связку двух осей.

\section{Статус машинной верификации в Lean~4.36.0 и честные границы}

Таблица~\ref{tab:status} разделяет результаты, уже проверенные ядром Lean~4.36.0 (\texttt{exit code: 0}, \texttt{0 sorryAx}), и компоненты, остающиеся на уровне спецификации, кандидата или открытой эмпирической задачи.

\begin{table}[!htbp]
\centering
\caption{Полная карта верифицированных модулей архитектуры v2.4 (Lean~4.36.0, 0 \texttt{sorryAx}).}
\label{tab:status}
\footnotesize
\begin{tabularx}{\linewidth}{@{}p{26mm}>{\hsize=1.25\hsize\raggedright\arraybackslash}X>{\hsize=0.85\hsize\raggedright\arraybackslash}X>{\hsize=0.90\hsize\raggedright\arraybackslash}X@{}}
\toprule
\textbf{Слой / модуль} & \textbf{Ключевые теоремы} & \textbf{Аксиомы Lean~4.36.0} & \textbf{Фактический статус} \\
\midrule
Ядро ревизии v2.4 & 15~деклараций (\texttt{pointwise\_\allowbreak unproved\_\allowbreak blocks\_\allowbreak any\_\allowbreak total}, \texttt{stepBranch\_\allowbreak fails\_\allowbreak without\_\allowbreak injectiveBin}, \texttt{metaFreeze\_\allowbreak preserves\_\allowbreak candidates}, \texttt{conditional\_\allowbreak two\_\allowbreak axis\_\allowbreak routing...} и др.) & 9 без аксиом; 2 \texttt{propext}; 1 \texttt{Quot.sound}; 3 стандартные & \textbf{Верифицировано}~\cite{strand2026kernel} (\texttt{1bb76af2...}, 605~стр., 0~warn) \\
Gate~11--12 v12.3 & 12~теорем (\texttt{checkMNCert\_\allowbreak sound}, \texttt{check1DCACert\_\allowbreak sound}, \texttt{checkFullChomskyProtocol\_\allowbreak sound}) & \texttt{propext}, \texttt{Classical.choice}, \texttt{Quot.sound} & \textbf{Верифицировано}~\cite{strand2026hl} (\texttt{763f59d7...}) \\
Gate~14 v2.4 (Порт) & 10~деклараций (\texttt{gaugeCert\_\allowbreak blocks\_\allowbreak single\_\allowbreak port\_\allowbreak id}, \texttt{tapPort\_\allowbreak breaks\_\allowbreak gaugeCert}, \texttt{spurious\_\allowbreak extension\_\allowbreak cannot\_\allowbreak repair\_\allowbreak refutation}) & 8 без аксиом; 2 \texttt{propext}, \texttt{Quot.sound} & \textbf{Верифицировано}~\cite{strand2026g14} (\texttt{34397f38...}, 272~стр., 0~warn) \\
Gate~15 v2.4 (Операторы) & 8~деклараций (\texttt{checkNonCommProtocol\_\allowbreak sound}, \texttt{noncomm\_\allowbreak refutes\_\allowbreak joint\_\allowbreak invariance}, анти-утечка \texttt{Held}/\texttt{N}) & 5 без аксиом; 3 \texttt{propext} & \textbf{Верифицировано}~\cite{strand2026g15} (\texttt{67fbae74...}, 231~стр., 0~warn) \\
Gate~16 v2.4 (Разрыв $\Delta_\varepsilon$) & 11~деклараций (\texttt{markov\_\allowbreak ineq}, \texttt{evalue\_\allowbreak false\_\allowbreak rejection\_\allowbreak bound}, \texttt{indep\_\allowbreak two\_\allowbreak step\_\allowbreak evalue\_\allowbreak product\_\allowbreak valid}, \texttt{den\_\allowbreak zero\_\allowbreak rejected\_\allowbreak by\_\allowbreak HasGapBound}) & 1 без аксиом; 3 \texttt{propext}; 7 \texttt{propext}, \texttt{Quot.sound} & \textbf{Верифицировано}~\cite{strand2026g16} (\texttt{9cf958f2...}, 304~стр., 0~warn) \\
Gate~13 (Уровень~4) & \texttt{QueueProgram}, $\forall m\,\neg\mathrm{StackProgram}(m)$ & --- & Кандидат (ожидает прогона Lean) \\
Эмпирика переноса & Выигрыш \texttt{StepBranch} и \texttt{meta-synset} по полной стоимости & --- & Открытая гипотеза (стенд в плане) \\
\bottomrule
\end{tabularx}
\end{table}

\begin{remark}[Точные границы применимости v2.4 по итогам формального аудита]
\begin{enumerate}
  \item \emph{Границы Gate~14:} \texttt{PortRefinement} формализует нестрогое уточнение ($\mathrm{proj}\circ P_2 = P_1$), тогда как строгость эскалации обеспечивается отдельно в \texttt{tapPort\_\allowbreak breaks\_\allowbreak gaugeCert}; условие \texttt{sameObs} в \texttt{SpuriousExtension} подаётся как гипотеза на всех $f:Q\to O$, а не выводится синтаксически из формы грамматики.
  \item \emph{Границы Gate~15:} ортогональность маски $K$ проверяется локально в заданной точке $x_0$, а теоремы смешения слоёв устанавливают существование двух артефактов: вводящей в заблуждение зависимости от порядка при нарушении \texttt{Held} и ложного семантического сдвига при невыровненном $N$.
  \item \emph{Границы Gate~16 v2.4:} требование $0 < \mathrm{den}$ включено в определение \texttt{HasGapBound} и подтверждено теоремой \texttt{den\_\allowbreak zero\_\allowbreak rejected\_\allowbreak by\_\allowbreak HasGapBound}; для вероятностной интерпретации оценок требуется дополнительное условие $\mathrm{totalMass}\,\mu>0$ (в текущем определении нулевая мера не исключена); правило \texttt{union\_\allowbreak gap\_\allowbreak bound} требует общего знаменателя $q$; \texttt{indep\_\allowbreak two\_\allowbreak step\_\allowbreak evalue\_\allowbreak product\_\allowbreak valid} доказывает произведение фиксированной пары шагов, а не произвольную остановку по данным (теорема Вилля).
  \item \emph{Границы ядра v2.4:} все 4 модуля компилируются с \texttt{0~warnings} в едином журнале \texttt{ALL\_\allowbreak 4\_\allowbreak MODULES\_\allowbreak v2\_\allowbreak 4\_\allowbreak VERIFICATION.log}; при этом \texttt{conditional\_\allowbreak two\_\allowbreak axis\_\allowbreak routing...} остаётся условной развилкой с постулированным \texttt{disjointCause} над свободными предикатами \texttt{axis1Refuted}/\texttt{axis2Gauge} (при ложных обоих предикатах ветка по умолчанию возвращает Ось~2, так как отдельный исход \texttt{InsufficientEvidence} отнесён к единому контроллеру следующей версии); \texttt{metaFreeze} переносит список кандидатов и префикс журнала, но пока не несёт полный сертификат разрыва и происхождения; в \texttt{gate12\_\allowbreak 13\_\allowbreak parametric\_\allowbreak lift} предикат \texttt{hPart} принимает посылку \texttt{True ->} на всём $Q_{\mathrm{mod}}$, а в \texttt{pointwise\_\allowbreak unproved\_\allowbreak blocks\_\allowbreak any\_\allowbreak total} несуществование доказано для тоталов с точным равенством области \texttt{tot.R = R}.
\end{enumerate}
\end{remark}

\section{Заключение}

В версии v2.4 дуальная пара \texttt{meta-synset}\,$\leftrightarrow$\,\texttt{invariant} охватывает обе оси исследовательских эскалаций: 44 машинно проверенные декларации в Lean~4.36.0 покрывают ядро ревизии v2.4, эскалацию наблюдательного порта (Gate~14), алгебру вмешательств с вычислимой маской $K$ (Gate~15) и количественный остаточный разрыв с точными e-value оценками (Gate~16). Центральными объектами стали запечатанное свидетельство \texttt{DirectEvidence s A c}, обязательный остаточный разрыв $\Delta(T)$ с доказательством разбиения \texttt{partition}, разведение моста полноты \texttt{bridgeComplete} и межобластного подъёма \texttt{liftScope}, зависимая лицензия вывода $\Gamma(\kappa)$ и ветвящийся выбор в меню $\calR$ с несущими предпосылками \texttt{SymbolLift}. Два непосредственных следующих шага --- локальная верификация Уровня~4 (Gate~13, \texttt{QueueProgram}) и сравнительный стенд полной стоимости ветвящегося контроллера \texttt{StepBranch} против линейной эскалации на отложенном наборе задач.

\begin{thebibliography}{99}\small
\bibitem{strand2026ms2}
А.\,Странд. \emph{Мета-синсет II: удостоверенная ревизия языка гипотез. Типизированный перенос, область вывода и проверяемая смена представления}, v0.3, рукопись, 2026.

\bibitem{strand2026hl}
А.\,Странд. \emph{Языки гипотез и доказательная эскалация представления. От конечных сертификатов невозможности к параметрическим теоремам о счётчиковых и стековых автоматах}, Gate~11--12 v12.3, 2026.

\bibitem{strand2026kernel}
А.\,Странд. \emph{epistemic\_\allowbreak revision\_\allowbreak kernel\_\allowbreak v2\_\allowbreak 4.lean}: формальное ядро системы ревизии v2.4 (Lean~4.36.0, SHA-256 \texttt{1bb76af2f6cc0334\allowbreak 8fa35e605d6ce41dfbd35cfcb9f5ae2a0ec8491c792319c7}, 15~деклараций, 0~warnings), 2026.

\bibitem{strand2026g14}
А.\,Странд. \emph{gate14\_\allowbreak port\_\allowbreak escalation\_\allowbreak v2\_\allowbreak 4.lean}: эскалация наблюдательного порта и запрет тупиковых расширений (Lean~4.36.0, SHA-256 \texttt{34397f3839bcc9cd\allowbreak b06335b1ca959e2f983dbad1635edf7430cb7d767e0ddcc6}, 10~деклараций, 0~warnings), 2026.

\bibitem{strand2026g15}
А.\,Странд. \emph{gate15\_\allowbreak intervention\_\allowbreak algebra\_\allowbreak v2\_\allowbreak 4.lean}: алгебра вмешательств, локальная вычислимая маска $K$ и коммутатор $[T_L,T_R]$ (Lean~4.36.0, SHA-256 \texttt{67fbae749cc0ab16\allowbreak 3468879396bca8e6f14f38f881ef63b78119888c4bad9c22}, 8~деклараций, 0~warnings), 2026.

\bibitem{strand2026g16}
А.\,Странд. \emph{gate16\_\allowbreak stochastic\_\allowbreak gap\_\allowbreak evalue\_\allowbreak v2\_\allowbreak 4.lean}: количественный остаточный разрыв $\Delta_\varepsilon$ ($0<\mathrm{den}$) и произведение двух независимых e-value (Lean~4.36.0, SHA-256 \texttt{9cf958f26f67e411\allowbreak 2367090ff5fbbcd895dcf9db849886494164a2b9fc640d99}, 11~деклараций, 0~warnings), 2026.

\bibitem{gentner1983}
D.\,Gentner. Structure-Mapping: A Theoretical Framework for Analogy. \emph{Cognitive Science}, 7:155--170, 1983.

\bibitem{goguen1992}
J.\,A.\,Goguen, R.\,M.\,Burstall. Institutions: Abstract Model Theory for Specification and Programming. \emph{Journal of the ACM}, 39(1):95--146, 1992.

\bibitem{blackwell1953}
D.\,Blackwell. Equivalent Comparisons of Experiments. \emph{Annals of Mathematical Statistics}, 24(2):265--272, 1953.

\bibitem{angluin1987}
D.\,Angluin. Learning Regular Sets from Queries and Counterexamples. \emph{Information and Computation}, 75(2):87--106, 1987.

\bibitem{myhill1957}
J.\,Myhill. Finite Automata and the Representation of Events. \emph{WADC Technical Report}, 57--624, 1957.

\bibitem{nerode1958}
A.\,Nerode. Linear Automaton Transformations. \emph{Proceedings of the American Mathematical Society}, 9(4):541--544, 1958.

\bibitem{lakatos1976}
I.\,Lakatos. \emph{Proofs and Refutations: The Logic of Mathematical Discovery}. Cambridge University Press, 1976.

\bibitem{huffman2013}
B.\,Huffman, O.\,Kun\v{c}ar. Lifting and Transfer: A Modular Design for Quotients. \emph{CPP 2013}, LNCS 8307, pp.\,131--146, 2013.

\bibitem{jha2017}
S.\,Jha, S.\,A.\,Seshia. A Theory of Formal Synthesis via Inductive Learning. \emph{Acta Informatica}, 54:693--726, 2017.
\end{thebibliography}

\end{document}
